\section{Conjugate Bayes, exponential families, Basu}
\label{sec:bayes}

\paragraph{Why these still need a script.}
Reviewed Bayes problems used an improper flat prior, a lognormal prior (MAP), and a two-point prior under absolute error. They never asked you to recognise a conjugate pair and write the posterior in closed form. Completeness has been used for UMVUE; ancillary statistics and Basu's theorem have never been named. Exponential families have been used implicitly; the natural-form definition has not been demanded.

\subsection{Lineage}

Conjugacy is the computational observation that, for an exponential family, a prior in the same family stays in that family \cite{gelman13,hogg19}. The decision-theoretic layer is: under \(L_{2}\) the Bayes estimator is the posterior mean; under \(L_{1}\) it is a posterior median. Basu (1955) is the independence of a complete sufficient statistic and an ancillary statistic; it is the usual way to see \(\bar X\perp S^{2}\) in a normal sample without computing a joint density. The natural exponential family packages sufficiency, completeness, and the Cram\'er--Rao variance in one Hessian \cite{casella02}.

\subsection{Concepts that belong on a script}

\begin{definition}[Natural exponential family]
\label{def:nef}
A density (or mass function) of the form
\[
f_{\eta}(x)=h(x)\exp\bigl(\eta^{\top}T(x)-A(\eta)\bigr)
\]
is a natural exponential family with natural parameter \(\eta\) and sufficient statistic \(T\). The cumulant function satisfies \(\nabla A=\mathbb{E}_{\eta}[T]\) and \(\nabla^{2}A=\Var_{\eta}(T)\). If the natural parameter space has nonempty interior, \(T\) is complete (and minimal sufficient). The Cram\'er--Rao bound for unbiased estimation of \(\nabla A(\eta)\) is \(\nabla^{2}A(\eta)/n\).
\end{definition}

\begin{definition}[Conjugate prior]
\label{def:conj}
A family \(\Pi\) of priors is conjugate for a sampling model if \(\pi\in\Pi\) implies that the posterior \(\pi(\,\cdot\mid x)\) stays in \(\Pi\). The three writeable pairs, with the posterior written out.

\begin{center}
\begin{tabular}{@{}lll@{}}
\toprule
Sampling & Prior & Posterior \\
\midrule
\(\mathrm{Binomial}(n,p)\) & \(\mathrm{Beta}(\alpha,\beta)\) & \(\mathrm{Beta}(\alpha+x,\,\beta+n-x)\) \\
\(\mathrm{Poisson}(\lambda)\) & \(\mathrm{Gamma}(\alpha,\beta)\) & \(\mathrm{Gamma}(\alpha+\sum x_{i},\,\beta+n)\) \\
\(N(\mu,\sigma^{2})\), \(\sigma^{2}\) known & \(N(\mu_{0},\tau^{2})\) & precision-weighted Gaussian \\
\bottomrule
\end{tabular}
\end{center}

The normal--normal update is
\[
\frac{\mu_{n}}{\tau_{n}^{2}}
=
\frac{\mu_{0}}{\tau^{2}}+\frac{n\bar X}{\sigma^{2}},
\qquad
\frac{1}{\tau_{n}^{2}}
=
\frac{1}{\tau^{2}}+\frac{n}{\sigma^{2}}.
\]
The \(L_{2}\) Bayes estimator is the posterior mean
\[
\mu_{n}
=
\frac{\tau^{-2}\mu_{0}+n\sigma^{-2}\bar X}{\tau^{-2}+n\sigma^{-2}},
\]
a convex combination of the prior mean and the MLE. The \(L_{1}\) Bayes estimator is a posterior median. For \(\mathrm{Beta}(\alpha',\beta')\) that median has no elementary closed form; the mean \(\alpha'/(\alpha'+\beta')\) is the \(L_{2}\) answer. The Poisson--Gamma posterior predictive is negative binomial: a new \(X_{n+1}\) has
\[
\mathbb{P}(X_{n+1}=k\mid\text{data})
=
\binom{k+\alpha'-1}{k}
\Bigl(\frac{\beta'}{\beta'+1}\Bigr)^{\alpha'}
\Bigl(\frac{1}{\beta'+1}\Bigr)^{k}
\]
in the shape--rate convention \(\mathrm{Gamma}(\alpha',\beta')\) with mean \(\alpha'/\beta'\) \cite{gelman13}.
\end{definition}

2026 Spring P9 already used the \(L_{1}\) rule on a two-point parameter space. The missing calculation is the Beta--Binomial (or Gamma--Poisson) update written in one line, plus the statement that the posterior mean is consistent if the prior charge of a neighbourhood of the true \(p\) is positive.

\begin{theorem}[Basu]
\label{thm:basu}
If \(T\) is a complete sufficient statistic for \(\theta\) and \(A\) is ancillary (its law does not depend on \(\theta\)), then \(T\) and \(A\) are independent for every \(\theta\).
\end{theorem}

The usual application: in \(N(\mu,\sigma^{2})\) with both parameters unknown, \((\bar X,S^{2})\) is complete sufficient and \(\bar X/S\) is not the example --- rather, for \(\mu\) unknown and \(\sigma^{2}\) known, \(S^{2}\) is ancillary for \(\mu\) after a location family argument, and Basu gives \(\bar X\perp S^{2}\). The 2024 Fall P10 paper asks that independence by a direct computation (mgf or rotational invariance). Basu is the named theorem that makes the computation a citation.

An ancillary example that belongs on a script: in a location family \(f(x-\theta)\), the spacings \(X_{(n)}-X_{(1)}\) (or the configuration \(X_{i}-\bar X\)) are ancillary.

\subsection{Connections}

\begin{itemize}
    \item Improper \(\pi(\theta)=1\) on \(U(0,\theta)\) (2025 Spring P7) is not conjugate. The posterior is proper for \(n\ge 2\); consistency is a separate argument.
    \item A lognormal prior MAP (2025 Fall P10) is a mode, not a posterior mean. Do not call it the Bayes estimator under \(L_{2}\).
    \item UMVUE arguments already use completeness. Basu is completeness plus ancillarity, used for independence, not for unbiased estimation.
    \item Fisher information of an exponential family is \(\nabla^{2}A\). Fisher information of a Gaussian linear model is \(Z^{\top}Z/\sigma^{2}\) (\cref{sec:lm}); same name, different Gram matrix.
\end{itemize}

\subsection{How it is examined}

A conjugate problem: name the pair, update the hyperparameters, write the posterior mean (and, if asked, the posterior median). A Basu problem: name \(T\) and \(A\), check completeness and ancillarity, conclude independence. An exponential-family problem: write \((\eta,T,A)\), read sufficiency and the CR bound from \(A\).

\subsection{Possible questions}

\begin{itemize}
    \item \(X_{i}\sim\mathrm{Bernoulli}(p)\), \(\mathrm{Beta}(\alpha,\beta)\) prior. Write the posterior and the \(L_{2}\) / \(L_{1}\) Bayes estimators.
    \item Poisson with a Gamma prior; predict a new \(X_{n+1}\) (posterior predictive is negative binomial).
    \item Prove \(\bar X\perp S^{2}\) in a normal sample by Basu, then by a density argument.
    \item Write the natural form of \(\mathrm{Gamma}(\alpha,\beta)\) with \(\alpha\) known. Identify \(T\) and \(A''\).
\end{itemize}

\subsection{Anchor problems}

{\footnotesize
\begin{center}
\begin{tabular}{@{}L{2.7cm}L{5.4cm}L{6.1cm}@{}}
\toprule
Anchor & What the paper wants & Near-miss \\
\midrule
(no paper) & conjugate update + posterior mean / median & improper / lognormal / two-point \\
(no paper) & Basu named & 2025 Spring P9 asked completeness only \\
(no paper) & natural exponential form & implicit in UMVUE papers \\
\done{2024 Fall P10} & \(\bar X\perp S^{2}\); UMVUE of \(\mu/\sigma\); admissibility of \(\bar X\) & draft in \texttt{p10-normal-mean-var.tex} \\
\bottomrule
\end{tabular}
\end{center}
}
